\section{The Propers: Detailed Commentary}
\zlabel{triptych:concise:commentary:start}
\subsection{What the servants know, and what they cannot know}
The servants know that the owner sowed good seed and that tares have
appeared. His reply confirms an enemy's action. Their recognition of a
real wrong therefore survives his refusal of their remedy. The reason
he gives concerns what uprooting would do to the wheat. Knowledge of
the wrong's origin does not itself establish that their proposed action
will preserve the good crop.

Chrysostom first explains the appearance of error. Heretics resemble the
truth they imitate, just as tares resemble wheat. False prophets follow
prophets and false apostles follow apostles. The order matters: good
has been sown before the enemy corrupts it. Chrysostom's account of
the servants' concern also matters. He recognises affection for the
crop in their proposal. Their fault concerns discretion about the
means of helping it. A good intention can accompany a dangerous
remedy.\footnote{Chrysostom, \work{Homilies on Matthew}, XLVI.1.}

His two explanations of injury reach different people. In the first,
violence against heretics produces war in which saints can die with
the guilty. Here the wheat already exists, and the means of removing
evil also destroy good people. In the second, a person now among the
tares can become wheat. Killing forecloses that person's repentance.
Here the endangered good is a possible conversion. Judging present
error accurately still leaves the person's future unknown. These
reasons explain why Chrysostom can oppose killing while continuing
to oppose heresy and permitting the suppression of its speech and
assemblies. His application distinguishes actions, rather than
pronouncing all forms of restraint equivalent.

Augustine develops the possibility of change by examining how a
parable signifies. Actual tares do not turn into a different botanical
species. Human beings, however, can repent. He compares the wicked
represented by tares with the hard ground of the preceding parable.
The hearer can be told both to have his hard ground ploughed and to
become wheat. Different images describe a person who can hear and
obey. Treating the image as a fixed human destiny would defeat the
sermon's appeal to change.\footnote{Augustine, Sermon XXIII
[Benedictine LXXIII], §§1--4.}

Augustine also distinguishes human servants from angelic reapers.
The angels identify correctly those whom they separate at the end;
human beings can take wheat for tares. This risk differs from the
risk of killing someone who would later convert. In one case the
present judgement is mistaken; in the other it is the future that
exceeds the judge's knowledge. Both restrict the claim to know whom
destructive action can safely remove. His inclusion of clergy and
laity prevents a hearer from finding an exempt class by looking at
a person's place in the assembly.

The distinction between people also bears on the speaker's view of
himself. Augustine addresses those he calls good with a warning about
falling. His concern extends beyond an error in classifying a neighbour:
a hearer can presume that his own perseverance is already secured.
The Collect's continuing request for protection makes sense for someone
in precisely this position. The Secret's direction of wavering hearts
concerns the offerers after pardon as well as before it. Mercy toward
the neighbour can be exercised by a person who still asks mercy for
himself. Augustine's call to imitate the good and reach God through
mercy addresses the whole assembly, across its differences of office
and present conduct.

The future fire and barn remain in both Fathers' accounts. Chrysostom
expects the unchanged tares to be burned; Augustine expects the
present mixture to end. Their hope for conversion concerns the time
before that separation. In the anagogical sense, the parable promises
an end that human servants cannot produce by hurrying it. The moral
question is what they do with the person still before them.

\subsection{How forgiveness can include correction}
Colossians places forbearance and admonition in the same life. Paul
first tells those who have received mercy to forgive a grievance.
He later commands teaching and admonishing one another. The injured
person therefore has more to do than refrain from retaliation.
Christ's word is to dwell abundantly among people who still need
to learn how to live together.

Chrysostom explains forgiveness from its stated comparison: the Lord
has forgiven. The first injury may have been entirely the other
person's fault, and the wrong may be serious. Neither fact cancels
the example Paul supplies. The forgiven person's response concerns
what his Lord has done for him, as well as what the offender has
done against him. Thanksgiving has a practical consequence here:
acknowledging mercy received changes how one treats another
servant.\footnote{Chrysostom, \work{Homilies on Colossians}, VIII,
on 3:12--15.}

His account of charity explains why an outward concession may remain
incomplete. A person can yield a claim while retaining coldness toward
the neighbour. Paul calls charity a bond because the virtues need
to remain joined in their relation to that person. Mercy as an inward
affection, rather than a bare concession, concerns what the injured
person actually desires for him. That is pertinent to the Gospel's
possible future wheat: desiring another's conversion requires him
to remain someone for whom one seeks a good.

The peace-as-umpire image gives a more exact account of this inward
choice. Revenge and endurance contend in the person's thoughts.
Peace awards the contest to endurance because God has reconciled
those whom Paul addresses. The image does not deny that a grievance
has arisen. It assigns the decision about response to a peace
already received from God. Chrysostom's textual lemma says ``peace
of God,'' while the appointed Latin says ``peace of Christ.''

This explains how the first interpretation can combine the Gospel
with Paul's instruction. The danger of uprooting concerns a particular
destructive act. Teaching and admonition aim at a person's change.
Chrysostom's warning about negligence in the field likewise demands
vigilance against error. The interval before harvest is occupied
by instruction and forgiveness, each governed by concern for the
person whom hasty action might injure.

Paul's order of teaching and thanksgiving gives correction a setting
within common worship. The people taught by Christ's word also sing
to God from their hearts. Chrysostom explains Scripture's abundance
as a lasting resource for people with household responsibilities.
A lesson learned from it remains available when a first moment of
grief has passed. The neighbour thus receives more than a verdict
on the present offence: instruction can prepare a different future
response. This is the kind of positive care available during the
Gospel's interval, within Paul's express commands to teach and admonish.

The same conduct enters the second interpretation through Mark's
account of prayer. Jesus's command to forgive follows his assurance
about receiving (Mark 11:22--26). Bede explains charitable faith at
that point, then directs the less perfect toward pardon and eternal
life while they forgive their offenders.\footnote{Bede, \work{Exposition
of Mark}, on 11:22--26, paragraphs beginning \textit{Et respondens Jesus}
and \textit{Et cum stabitis ad orandum}.} In relation to the Gospel,
forgiveness concerns the neighbour's possible future; in relation
to the Communion, it concerns the person approaching God. Those
consequences are compatible, but they answer distinct questions.

\subsection{Why a promised gift is still requested}
Jeremias's letter directs the exiles toward a settled obedience.
They must build and plant where they are, despite rival promises
of a speedy return. Jerome's explanation of God's invitation to
pray begins within this conflict over expectation. Confidence in
God concerns the word he has actually spoken, including the life
he commands before restoration. The desired quick return cannot
be made true by the strength of the captives' wish.

Jerome then distinguishes God's power from the people's commanded
response. God could restore them without their prayer. Yet his
command to call and seek makes prayer part of their reception of
the promise. Divine sufficiency therefore supplies no reason for
disobedience. The whole-hearted seeker answers the giver who has
already addressed him.\footnote{Jerome, \work{Commentary on Jeremias},
V, on 29:10--14, PL 24, cols.~893--894.}

Augustine reaches a related difficulty from Psalm 84. The psalm
thanks God for turning captivity and still asks him to turn the
people. Augustine explains how the prophet can see the promised
restoration as accomplished and pray for its coming. The certainty
and the request concern the same gift. In his Christian exposition,
forgiveness releases the captivity of sin, so the petition concerns
the people's own turning as well as a change in circumstances.
The literal exiles addressed by Jeremias remain distinct from this
Christian explanation of the psalm.\footnote{Augustine,
\work{Exposition of Psalm 84}, §§2--6.}

The Gradual presents thanksgiving in another demanding context.
Psalm 43 remembers God's victories and then laments defeat despite
the people's fidelity. Augustine hears the martyrs through Paul's
use of the psalm's sheep-for-slaughter words in Romans 8. He also
explains the past language of rescue by prophetic certainty,
comparing other prophetic past tenses and the psalm's future praise.
The interpretation allows assured praise while deliverance remains
awaited. It leaves present suffering fully real.\footnote{Augustine,
\work{Exposition of Psalm 43}, §§1,3--8.}

Jerome's further explanation of return concerns final rest beyond
the world's confusion. His anagogical extension starts from the
exiles' historical restoration. It gives the second interpretation
a final object of hope beyond relief from a present trial. Augustine's
Christian account of restored captivity likewise concerns the people
released from sin. Both explanations attend to what happens to the
recipients. Their return includes the turning by which they receive
mercy, rather than leaving their desires and conduct untouched.

These explanations yield different consequences for the two
interpretations. Patience toward the neighbour can rest on the
certainty that judgement belongs to God: private vengeance need
not make the promised end visible now. Persistent prayer can rest
on the same God's promise while accepting an interval in which
outward circumstances remain difficult. Neither inference identifies
the exiles' neighbours or the psalm's enemies with particular tares
in the Gospel. The comparison concerns obedience and expectation
during a present whose end God has promised.

\subsection{What confidence asks from mercy}
Augustine begins Psalm 129 with Jonas praying inside the fish.
The image establishes that no physical depth can keep the cry
from God. He then examines a more dangerous depth: the condition
of a sinner who no longer asks for rescue. Continued prosperity
can be misread as evidence that sin has no consequence. The
petitioner who cries has begun to reject that self-satisfaction.
Acknowledging need and expecting mercy already belong to hope.
\footnote{Augustine, \work{Exposition of Psalm 129}, §§1--2.}

The psalm asks who could stand if the Lord marked iniquities
strictly. Augustine attends to the question's reach: the speaker
does not exempt everyone else while confessing his own weakness.
Confidence must therefore rest in God's propitiation. Bellarmine
adds a distinction within the opening cry. Depth can name misery
and also the heart's inward attention. An external calamity does
not itself make the person attentive to God; serious prayer gives
the heart to its asking.\footnote{Bellarmine, \work{Commentary on
the Book of Psalms}, Ps 129, explanations of verses 1--2,
O'Sullivan's abridged translation.}

The repeated cry in Alleluia and Offertory can thus be understood
as hopeful penitence. At the offering, it is followed by the
Secret's request for pardon and the direction of wavering hearts.
The saving good includes what the petitioner will become. This
relation helps specify the confidence with which the Communion's
promise is heard. Bede's account of faith working through charity
already includes forgiveness of an actual offender.

Aquinas gives a further distinction in his doctrinal exposition
of prayer. Merit concerns the charity from which prayer proceeds;
obtaining the requested good, or impetration, concerns God's gracious
gift. A devout act can remain good when a particular temporal
request would not conduce to the person's salvation. God knows
whether that outcome would help the person reach beatitude.
\footnote{Aquinas, \work{Summa theologiae} II--II, q.~83, a.~15,
corpus and replies 1--3. This is doctrinal illumination of the
Communion's question, not direct exegesis of Mark 11.}

His four conditions specify asking for oneself, for what is necessary
to salvation, piously and perseveringly. A useful gift can also be
delayed until a fitting time. Consequently an ungranted temporal
petition cannot simply be diagnosed as deficient faith. The distinction
concerns both the desired good and the way God's mercy grants it.
It gives precision to the second interpretation: the promise warrants
persevering petition for salvation without making the present wish
the measure of God's gift.

Aquinas also distinguishes a sinner's evil request from a good
request made by a person who has sinned. God loves that person's
nature while hating the sin and can hear him out of mercy.\footnote{Aquinas,
\work{Summa theologiae} II--II, q.~83, a.~16, corpus.}
The petitioner from the depths can ask for pardon before becoming
innocent of every fault. In the first interpretation this dependence
checks the corrector's presumption; in the second it gives the
penitent reason to approach God now.

\subsection{The salvation begun, and its completion}
Augustine explains the psalm's propitiation by Christ's sacrifice:
innocent blood redeems the guilty. He then joins forgiveness of the
past with warning concerning the future. That sequence explains
why the Secret asks both absolution and direction. The offered
sacrifice and the petition belong to the prayer itself; Augustine
explains the saving action of Christ at the psalm's words.
\footnote{Augustine, \work{Exposition of Psalm 129}, §3.}

His account of waiting proceeds to the morning watch and Christ's
resurrection. The reason a member can hope from the risen Head
is the flesh Christ took from humanity. He received what he would
offer for it; that flesh has risen as first-fruits. The body has
therefore been given a ground for expecting its own resurrection,
although its members have not yet undergone it. Redemption from
iniquities addresses their sin within that hope.\footnote{Augustine,
\work{Exposition of Psalm 129}, §§4--8.}

The Postcommunion makes its own distinction between the pledge
received through the mysteries and salvation's effect still sought.
Beside Augustine's exposition, reception can be understood as a
gift that awakens desire for completion. This gives the second
interpretation its particular allegorical and anagogical account:
Christ's sacrifice makes the sinner's approach possible, and his
resurrection grounds the body's final hope. The literal petition
asks to attain the effect of the salvation already pledged.

The first interpretation reaches the final end through the Gospel's
separation. Christ's gathering and judgement end the interval in
which people can still change. The sacramental pledge therefore
does not exempt the recipient from perseverance. Augustine warns
the apparently secure hearer to take care lest he fall. The
Postcommunion's continuing petition concerns people who have
already received a real gift and still need its completion.

During that interval, Colossians commands instruction, heartfelt
song and thanksgiving. Chrysostom addresses the abundance of
Scripture to people responsible for families as well as teachers.
Aquinas explains how doing all in Christ's name includes habitual
intention, without requiring a continuously explicit mental act.
Ordinary work can thus be directed to God by people still waiting
for salvation's effect.\footnote{Chrysostom, \work{Homilies on
Colossians}, IX, opening on 3:16--17; Aquinas, \work{Commentary on
Colossians}, chapter 3, lectio 3, on 3:16--17, Dessain II, p.~442.}
The Preface offers the thanksgiving to which Paul calls them: worship
of the Father with the Son and Holy Spirit, in distinct Persons
and one divine essence. Its confession concerns God; Paul's
command to one body concerns redeemed people learning charity.

The interpretations accordingly conclude with different consequences.
Care for the neighbour must reckon with mistaken identification,
harm to the innocent and a future conversion one cannot foresee.
Confident prayer must seek a saving good through mercy, even while
the desired temporal result remains absent. Both concern people
who forgive because they have been forgiven and ask because God
has already begun to give.
